\begin{thebibliography}{99}
\bibitem{OAI4}
OpenAI.
\emph{Every Four-Dimensional Kakeya Set Has Full Hausdorff Dimension}.
2026年9月24日预印本，175页；固定仓库版本
\href{https://github.com/openai/math/tree/1581d4cbdd69a0faf34b2089614fe2a814316a6a/preprints/Every-four-dimensional-Kakeya-set-has-full-Hausdorff-dimension-September-24-2026}{openai/math，提交1581d4cb}。

\bibitem{KT}
Nets Hawk Katz and Terence Tao.
\emph{New bounds on Kakeya problems}.
Journal d'Analyse Math\'ematique \textbf{87} (2002), 231--263.
本文所用预印本：\href{https://arxiv.org/abs/math/0102135v1}{arXiv:math/0102135v1}，2001；第5--6节，特别是定理6.2。

\bibitem{HRZ}
Jonathan Hickman, Keith M. Rogers and Ruixiang Zhang.
\emph{Improved bounds for the Kakeya maximal conjecture in higher dimensions}.
\href{https://arxiv.org/abs/1908.05589v1}{arXiv:1908.05589v1}，2019；特别是Figure 5。

\bibitem{GLX}
Chuanwei Gao, Diankun Liu and Yakun Xi.
\emph{Curved Kakeya sets and Nikodym problems on manifolds}.
\href{https://arxiv.org/abs/2503.11574v1}{arXiv:2503.11574v1}，2025；特别是Table 1的欧氏对照。

\bibitem{GWZ}
Larry Guth, Hong Wang and Joshua Zahl.
\emph{A streamlined proof of the Kakeya set conjecture in $\R^3$}.
\href{https://arxiv.org/abs/2601.14411v2}{arXiv:2601.14411v2}，2026；定理1.1。

\bibitem{WZ}
Hong Wang and Joshua Zahl.
\emph{Sticky Kakeya sets and the sticky Kakeya conjecture}.
\href{https://arxiv.org/abs/2210.09581v2}{arXiv:2210.09581v2}；定义4.2及定理5.2。

\bibitem{CV}
Anthony Carbery and Stef\'an Ingi Valdimarsson.
\emph{The endpoint multilinear Kakeya theorem via the Borsuk--Ulam theorem}.
\href{https://arxiv.org/abs/1205.6371v1}{arXiv:1205.6371v1}，2012；定理1。

\bibitem{OSS}
Tuomas Orponen, Nicolas de Saxc\'e and Pablo Shmerkin.
\emph{On the Fourier decay of multiplicative convolutions}.
\href{https://arxiv.org/abs/2309.03068v2}{arXiv:2309.03068v2}；定理1.5。

\bibitem{ArchiveSource}
本项目研究档案。
\emph{R0038：任意集合归约、原证明拼接与有限/Hausdorff出口}。
2026年10月17日；固定版本的
\href{https://github.com/lixiang90/kakeya_research/blob/774c3614ec3e810c127c875a9f5abb60748be848/notes/10_source_audit/source_assembly/R0038_reduction_source_assembly_and_finite_readout_full_proof.md}{RED-001完整证明}。
内部交叉审查，非外部同行评审。

\bibitem{ArchiveDIR}
本项目研究档案。
\emph{R0009：稀疏方向放大与加权方向化约}。
固定版本的
\href{https://github.com/lixiang90/kakeya_research/blob/774c3614ec3e810c127c875a9f5abb60748be848/notes/40_multiscale/plank_inputs/R0009_sparse_direction_amplification.md}{稀疏方向证明}及
\href{https://github.com/lixiang90/kakeya_research/blob/774c3614ec3e810c127c875a9f5abb60748be848/notes/40_multiscale/plank_inputs/R0009_weighted_direction_reduction.md}{权重化约}。

\bibitem{ArchiveFSD}
本项目研究档案。
\emph{R0029：任意环境的有限投影、原标记与六维种子}。
固定版本的
\href{https://github.com/lixiang90/kakeya_research/blob/774c3614ec3e810c127c875a9f5abb60748be848/notes/20_dimension_lifting/finite_seed_transfer/R0029_projection_and_prior_audit.md}{有限投影证明}与
\href{https://github.com/lixiang90/kakeya_research/blob/774c3614ec3e810c127c875a9f5abb60748be848/notes/20_dimension_lifting/finite_seed_transfer/R0029_conditional_seed_and_six_dimensional_bound.md}{种子及递推整合}。
R0029最初的紧集措辞已由R0038任意集合读出扩展。

\bibitem{ArchiveDIM}
本项目研究档案。
\emph{DIM-001：任意六维Kakeya集合的Hausdorff下界17/4}。
2026年10月19日；
\href{https://github.com/lixiang90/kakeya_research/blob/774c3614ec3e810c127c875a9f5abb60748be848/results/reviewed/DIM-001_six_dimensional_Hausdorff_research_report.md}{固定版本成果报告及审查定位}。
\end{thebibliography}
